\begingroup
\renewcommand{\arraystretch}{1.45}
\InterpolationArticle{19}
Besonders häufig und daher der größten Aufmerksamkeit würdig ist der Fall, in dem
die Werte des Bogens \(x\), für die Werte der Funktion \(X\) gegeben sind, eine arithmetische Folge bilden, deren Gliederdifferenz dem Quotienten entspricht, der
bei der Teilung des ganzen Umfangs in ebenso viele Teile entsteht, wie Glieder vorhanden sind. Würde also
die Folge über ihr letztes Glied hinaus fortgesetzt, so würden ihre ersten Glieder, um einen ganzen Umfang vermehrt, der Reihe nach wiederkehren. Nachdem wir oben die Hauptpunkte der allgemeinen Untersuchung abgeschlossen haben, werden wir diesen Fall, in dem
die gegebenen Glieder gleichsam eine vollständige \emph{Periode} bilden, als krönenden Abschluss gesondert ausführlicher behandeln. Bevor wir jedoch diese Untersuchung selbst beginnen, sind zwei
Hilfssätze voranzuschicken.

\textsc{Erster Hilfssatz.} \emph{Bilden die Bögen \(a,b,c,d\) usw., deren Anzahl \(\mu\) ist,
eine arithmetische Folge, in der die Gliederdifferenz \(b-a=c-b=d-c\)
usw. \(=\frac{360^{0}}{\mu}\) ist, so wird das Produkt aus \(\mu\) Faktoren \(P=\)}
\SourceLeaf{303}
\[
\begin{aligned}
&\sin\frac12(t-a)\sin\frac12(t-b)\sin\frac12(t-c)\sin\frac12(t-d)\ \mathrm{usw.}
\end{aligned}
\]
\noindent\emph{\(=\mp\frac{\sin\frac12\mu(t-a)}{2^{\mu-1}}\), wobei das obere oder das untere Zeichen gilt, je nachdem, ob \(\mu\) gerade oder ungerade ist.}

\emph{Beweis.} Dieser Satz lässt sich zwar leicht aus einem anderen herleiten, den der berühmte \textsc{Euler} in \emph{Introd. in Anal. Infin.} I, \S. 240 angegeben hat; damit hier jedoch nichts zu fehlen scheint, werden wir einen anderen Beweis kurz erläutern.

\Needspace{30\baselineskip}
I. Wird das Produkt \(P\) zunächst in die Form
\[
\begin{aligned}
&\cos\frac12(t-a-180^{0})\cos\frac12(t-b-180^{0})\\[5pt]&\qquad\cos\frac12(t-c-180^{0})\cos\frac12(t-d-180^{0})\mathinner{.\,.\,.}
\end{aligned}
\]
\noindent gebracht und dann nach den bekannten Regeln der Sinusrechnung in eine Summe von Kosinusgrößen entwickelt, so sieht man leicht, dass daraus eine Funktion entweder von der Form \(F\) oder von der Form \(G\)
(Artikel 13) hervorgeht, die bis auf \(\frac12\mu^{\mathrm{te}}\) Ordnung steigt, und dass ihr \emph{höchstes} Glied
\[
\begin{aligned}
&=\frac{1}{2^{\mu-1}}\cos\frac12(\mu t-a-b-c-d-\mathrm{usw.}-\mu\times180^{0})\\[5pt]
&=\frac{1}{2^{\mu-1}}\cos\bigl(\frac12\mu(t-a-360^{0})+90^{0}\bigr)\\[5pt]
&=-\frac{1}{2^{\mu-1}}\sin\frac12\mu(t-a-360^{0})\\[5pt]
&=-\frac{1}{2^{\mu-1}}\sin\frac12\mu(t-a)\cos\mu\times180^{0}\\[5pt]
&=\mp\frac{1}{2^{\mu-1}}\sin\frac12\mu(t-a)
\end{aligned}
\]
\noindent wird, wobei das obere Zeichen für einen geraden Wert von \(\mu\) gilt, das untere für einen ungeraden.

\Needspace{11\baselineskip}
II. Da \(\sin\frac12\mu(t-a)\) \(=0\) wird, wenn \(t\) einer der Werte \(a.\ b,c,d\) usw.
zugewiesen wird, ist nach dem Satz von Artikel 14 \(\sin\frac12\mu(t-a)\) durch das Produkt \(P\) teilbar;
der Quotient ist jedoch notwendig eine bestimmte Größe \(=h\), da \(\sin\frac12\mu(t-a)\) und \(P\)
von derselben Ordnung sind. Aus der Gleichung \(P=\frac1h\sin\frac12\mu(t-a)\) folgt aber offenbar (nach I),
dass \(h=\mp2^{\mu-1}\) ist, also \(P=\mp\frac{\sin\frac12\mu(t-a)}{2^{\mu-1}}\)\quad Q.~E.~D.

\Needspace{32\baselineskip}
\textsc{Zweiter Hilfssatz.} \emph{Verhalten sich die Bögen \(a,b,c,d\mathinner{.\,.\,.}\) so wie im ersten Hilfssatz, so wird das Produkt aus \(\mu\) Faktoren}
\[
\begin{aligned}
&(\cos t-\cos a)(\cos t-\cos b)(\cos t-\cos c)(\cos t-\cos d)\mathinner{.\,.\,.}
\end{aligned}
\]
\noindent\emph{\(=\frac{\cos\mu t-\cos\mu a}{2^{\mu-1}}\)}

\emph{Beweis.} Da
\SourceLeaf{304}
\[
\begin{aligned}
&\cos t-\cos a=2\sin\frac12(t-a)\sin\frac12(-t-a)\\[5pt]
&\cos t-\cos b=2\sin\frac12(t-b)\sin\frac12(-t-b)\\[5pt]
&\cos t-\cos c=2\sin\frac12(t-c)\sin\frac12(-t-c)\\[5pt]
&\cos t-\cos d=2\sin\frac12(t-d)\sin\frac12(-t-d)\\[5pt]
&\mathrm{usw.}
\end{aligned}
\]
\noindent wird das Produkt aus diesen Faktoren nach dem ersten Hilfssatz
\[
\begin{aligned}
&=2^{\mu}\times\mp\frac{\sin\frac12\mu(t-a)}{2^{\mu-1}}\times\mp\frac{\sin\frac12\mu(-t-a)}{2^{\mu-1}}\\[5pt]
&=\frac{\sin\frac12\mu(t-a)\sin\frac12\mu(-t-a)}{2^{\mu-2}}=\frac{\cos\mu t-\cos\mu a}{2^{\mu-1}}\quad\mathrm{Q.\ E.\ D.}
\end{aligned}
\]

\Needspace{19\baselineskip}
Da
\[
\begin{aligned}
&\sin\frac12\mu(t-a)=-\sin\frac12\mu(t-b)\\[5pt]&=\sin\frac12\mu(t-c)=-\sin\frac12\mu(t-d)\ \mathrm{usw.}
\end{aligned}
\]
\noindent und außerdem \(\cos\mu a=\cos\mu b=\cos\mu c=\cos\mu d\) usw. gilt, wird offenbar das Produkt
im ersten Hilfssatz auch
\[
\begin{aligned}
&=\pm\frac{\sin\frac12\mu(t-b)}{2^{\mu-1}}\\[5pt]&=\mp\frac{\sin\frac12\mu(t-c)}{2^{\mu-1}}=\pm\frac{\sin\frac12\mu(t-d)}{2^{\mu-1}}\ \mathrm{usw.}
\end{aligned}
\]
\noindent und ebenso das Produkt im zweiten Hilfssatz auch
\[
\begin{aligned}
&=\frac{\cos\mu t-\cos\mu b}{2^{\mu-1}}=\frac{\cos\mu t-\cos\mu c}{2^{\mu-1}}\\[5pt]&=\frac{\cos\mu t-\cos\mu d}{2^{\mu-1}}\ \mathrm{usw.}
\end{aligned}
\]

\endgroup
